% Shared preamble for the repository diagrams. Each figure is compiled twice by build.sh: once as is
% (light theme) and once with \darktheme defined, which swaps the palette for a dark background.
\documentclass[tikz,border=4pt]{standalone}
\usepackage[T1]{fontenc}
\usepackage{lmodern}
\renewcommand{\familydefault}{\sfdefault}
\usepackage{amsmath,amssymb}
\usepackage{sansmath}
\sansmath
% Narrow boxes read better with whole words than with hyphenated ones.
\hyphenpenalty=10000
\exhyphenpenalty=10000
\usetikzlibrary{arrows.meta,positioning,calc,shapes.misc}

% Fill colors encode the role of a step: amber = test or rule, teal = forward simulation,
% purple = optimization or candidate generation, green = state or returned action, red = adverse
% outcome, blue = other processing.
\ifdefined\darktheme
  \definecolor{ink}{HTML}{E6EDF3}
  \definecolor{edge}{HTML}{8B949E}
  \definecolor{blue}{HTML}{1B2E45}
  \definecolor{teal}{HTML}{153C39}
  \definecolor{green}{HTML}{1C3823}
  \definecolor{amber}{HTML}{3F3313}
  \definecolor{red}{HTML}{442022}
  \definecolor{purple}{HTML}{30264A}
  \definecolor{grey}{HTML}{262B30}
  \definecolor{yes}{HTML}{56D364}
  \definecolor{no}{HTML}{F47067}
\else
  \definecolor{ink}{HTML}{17202A}
  \definecolor{edge}{HTML}{5D6670}
  \definecolor{blue}{HTML}{DDEBF7}
  \definecolor{teal}{HTML}{DDF2EF}
  \definecolor{green}{HTML}{E4F2E5}
  \definecolor{amber}{HTML}{FFF0CC}
  \definecolor{red}{HTML}{F8DDDD}
  \definecolor{purple}{HTML}{EAE2F5}
  \definecolor{grey}{HTML}{ECEFF1}
  \definecolor{yes}{HTML}{236B3B}
  \definecolor{no}{HTML}{A33636}
\fi

\tikzset{
  every picture/.style={color=ink, font=\small},
  box/.style={draw=edge, rounded corners=3pt, align=center, inner sep=4pt,
              text width=#1, minimum height=12mm, text=ink},
  box/.default=2.7cm,
  proc/.style={box=#1, fill=blue},   proc/.default=2.7cm,
  sim/.style={box=#1, fill=teal},    sim/.default=2.7cm,
  opt/.style={box=#1, fill=purple},  opt/.default=2.7cm,
  out/.style={box=#1, fill=green},   out/.default=2.7cm,
  bad/.style={box=#1, fill=red},     bad/.default=2.7cm,
  rule/.style={box=#1, fill=amber},  rule/.default=2.7cm,
  ask/.style={box=#1, fill=amber, chamfered rectangle, chamfered rectangle xsep=2pt,
              rounded corners=0pt, line width=0.6pt},
  ask/.default=2.7cm,
  arr/.style={-{Latex[length=2mm]}, semithick, draw=edge},
  darr/.style={arr, dashed},
  lab/.style={font=\scriptsize\bfseries, inner sep=1.5pt},
  ylab/.style={lab, text=yes}, nlab/.style={lab, text=no},
  note/.style={font=\scriptsize, text=ink, align=center},
}

\begin{document}
% Map of the rules over the two loadings the controller reads: the current peak loading (x) and the
% peak loading forecast for do-nothing over the next two steps (y). Scales: x 15 cm per unit from
% 0.70, y 7 cm per unit from 0.80.
\begin{tikzpicture}
\def\X#1{(#1-0.70)*15}
\def\Y#1{(#1-0.80)*7}
% regimes of the reactive controller
\fill[green] (0,0) rectangle ({\X{0.90}},{\Y{1.60}});
\fill[amber] ({\X{0.90}},0) rectangle ({\X{1.40}},{\Y{1.30}});
\fill[red] ({\X{0.90}},{\Y{1.30}}) rectangle ({\X{1.40}},{\Y{1.60}});
\draw[draw=edge] (0,0) rectangle ({\X{1.40}},{\Y{1.60}});
\foreach \r in {0.95,1.00}{\draw[draw=edge, dashed] ({\X{\r}},0)--({\X{\r}},{\Y{1.60}});}
\node[note, text width=2.7cm] at ({\X{0.80}},{\Y{1.20}}) {\textbf{comfortable}\\$\rho_{\max}\le0.90$\\[4pt]lift curtailment, charge storage, preserve ramp margin, return redispatch toward its economic set point};
\node[note, text width=5.6cm] at ({\X{1.17}},{\Y{1.03}}) {\textbf{preventive}\\[2pt]DC optimal power flow if a line exceeds its limit; anticipatory dispatch against the likeliest outage; the same program again while an overflow counter runs; otherwise do nothing};
\node[note, text width=5.6cm] at ({\X{1.17}},{\Y{1.45}}) {\textbf{danger}: candidate search\\(also entered when an announced outage makes the do-nothing forecast end the episode)};
% axes
\foreach \r in {0.70,0.80,0.90,0.95,1.00,1.10,1.20,1.30,1.40}{\draw[draw=edge] ({\X{\r}},0)--++(0,-0.12) node[below, font=\scriptsize, text=ink]{\r};}
\foreach \f in {0.80,1.00,1.30,1.60}{\draw[draw=edge] (0,{\Y{\f}})--++(-0.12,0) node[left, font=\scriptsize, text=ink]{\f};}
\node[note, font=\small] at ({\X{1.05}},-0.95) {current peak loading $\rho_{\max}(o_t)$};
\node[note, font=\small, rotate=90] at (-1.25,{\Y{1.20}}) {forecast peak loading of do-nothing, $F_2$};
\node[note, anchor=south east] at ({\X{1.40}},{\Y{1.60}}) {$F_2=+\infty$ when the forecast ends the episode};
% where the other rules act, on the same horizontal scale
\def\rulebar#1#2#3{\draw[draw=hledge, fill=hl, rounded corners=1pt] ({\X{#1}},#3) rectangle ({\X{#2}},{#3+0.3});}
\rulebar{0.70}{0.95}{6.10}
\node[anchor=west, font=\scriptsize, text=ink] at ({\X{0.70}},6.25) {\textbf{2} reopens withheld lines};
\rulebar{0.70}{1.00}{6.55}
\node[anchor=west, font=\scriptsize, text=ink] at ({\X{1.00}+0.1},6.70) {\textbf{2} does not reconnect a region's only disconnected attackable line};
\rulebar{0.92}{1.00}{7.00}
\node[anchor=west, font=\scriptsize, text=ink] at ({\X{1.00}+0.1},7.15) {\textbf{3} guarded-line rule, for two lines};
\rulebar{0.95}{1.40}{7.45}
\node[anchor=west, font=\scriptsize, text=ink] at ({\X{0.95}+0.1},7.60) {\textbf{3} AC-anchored dispatch competes with the relief action};
\rulebar{1.00}{1.40}{7.90}
\node[anchor=west, font=\scriptsize, text=ink] at ({\X{1.00}+0.1},8.05) {\textbf{4} may replace the action of \textbf{3} (a line is overloaded)};
\lanemap{0.95,9.0}{3}{which rule acts at which loading}
\end{tikzpicture}
\end{document}
